\section{Conclusion}
\label{sec:conclusion}

We presented \ARES{}, a compact formal reliability framework for
detecting anomalous behaviour in agentic AI systems deployed in
cybersecurity environments. The conference contribution is narrow:
the SKC mismatch model produces a learnable ARS that generalises
across threat classes (test F1 0.96), and the Epistemic Gap
metric formalises the hallucination precondition as
\textit{confident ignorance} and decomposes false-negative risk into
knowledge and calibration components. A cascade reliability bound for
tightly-coupled pipelines and a temporal-drift model for knowledge
ageing are reported as supporting results.

These results suggest several practical implications: ARS supports
pre-deployment reliability screening, the Epistemic Gap supports
runtime hallucination-risk monitoring, the temporal-drift model
supports evidence-based knowledge maintenance, and the cascade bound
informs conservative multi-agent design. The framework provides an
initial formal foundation for trustworthy deployment of agentic AI in
cybersecurity, and may generalise to other domains where autonomous
agents make consequential decisions under incomplete knowledge.
Future work extends \ARES{} with real-world SOC data, cyclic
multi-agent architectures, and adversarial mismatch injection.
The full evaluation testbed, the labelled \SAFK{} corpus, and
extended materials (full proofs, per-class breakdowns, contingency
tables, and weight-sensitivity sweeps) are publicly
available.\footnote{\url{https://github.com/IntelEyes/ARES-Framework}}
